\chapter{Architecture Système, Implémentation et Validation de Navigation}
\label{chap:realization}

\section{Introduction}
Après avoir formalisé les exigences et la conception conceptuelle, ce chapitre expose la réalisation logicielle et l'assurance qualité intégrale du module \textbf{\odooModule}, développé pour \textbf{\companyName}. Nous détaillons l'architecture multi-niveaux sous Odoo 18, l'inventaire quantitatif du code, la spécification approfondie des modèles de données ORM, les algorithmes prédictifs statistiques et la couche d'intelligence artificielle conversationnelle BYO-LLM (Google Gemini et Cerebras). Nous consacrons une section majeure à la stratégie de **tests E2E de navigation et d'interface utilisateur sous Playwright**, mettant en lumière la validation rigoureuse des vues (Kanban, listes groupées, formulaires, calendrier, multi-sites, wizards modaux) et les dix anomalies d'ergonomie et de navigation résolues. Enfin, nous présentons la télémétrie complète attestant de la réussite de 138 tests backend, 46 tests mock LLM, tests live d'API et 32 tests de navigation Playwright.

\section{Architecture Système Multi-Niveaux et Couche BYO-LLM}
Le progiciel s'appuie sur une architecture logicielle modulaire en quatre couches, augmentée d'un banc de test automatisé de navigation de bout en bout, tel que schématisé dans la figure \ref{fig:odoo_architecture}.

\begin{figure}[!ht]
    \centering
    \begin{adjustbox}{max width=0.95\textwidth}
    \begin{tikzpicture}[
        layer/.style={rectangle, draw=odooDark, thick, rounded corners=6pt, minimum width=13.5cm, minimum height=1.4cm, align=center, font=\small},
        comp/.style={rectangle, draw=odooPurple, fill=lightPurple, rounded corners=4pt, font=\scriptsize\bfseries, minimum height=0.8cm, align=center}
    ]
        % Layers
        \node[layer, fill=blue!5] (pres) at (0, 4.5) {};
        \node[font=\bfseries\color{odooBlue}, anchor=west] at (-6.5, 4.9) {Couche Présentation (Odoo Web Client \& Validation Navigation)};
        \node[comp] (p1) at (-4.2, 4.2) {Vues XML Odoo 18\\(Card / List / Form)};
        \node[comp] (p2) at (-1.3, 4.2) {Framework OWL\\(Composants JS)};
        \node[comp] (p3) at (1.5, 4.2) {Générateur QR Code\\\& Rapports QWeb PDF};
        \node[comp] (p4) at (4.8, 4.2) {Banc Playwright E2E\\(Validation Navigation)};

        \node[layer, fill=purple!5] (biz) at (0, 2.2) {};
        \node[font=\bfseries\color{odooPurple}, anchor=west] at (-6.5, 2.6) {Couche Métier, Prédictif \& BYO-LLM (Python 3.12 / ORM Odoo 18)};
        \node[comp] (b1) at (-4.2, 1.9) {Modèles Métier\\\& Multi-Sites};
        \node[comp] (b2) at (-1.3, 1.9) {Moteurs d'États\\\& Workflows};
        \node[comp] (b3) at (1.5, 1.9) {Algorithmes Santé\\\& Risque MTBF};
        \node[comp] (b4) at (4.8, 1.9) {Couche BYO-LLM\\(Gemini / Cerebras)};

        \node[layer, fill=green!5] (data) at (0, -0.1) {};
        \node[font=\bfseries\color{greenOk}, anchor=west] at (-6.5, 0.3) {Couche Données \& Persistance (PostgreSQL 15+)};
        \node[comp] (d1) at (-4.0, -0.4) {Tables Relationnelles ORM};
        \node[comp] (d2) at (0.0, -0.4) {Contraintes Intégrité \& Triggers};
        \node[comp] (d3) at (4.2, -0.4) {Row-Level Security (ACL \& Rules)};

        \node[layer, fill=orange!5] (ext) at (0, -2.4) {};
        \node[font=\bfseries\color{orangeWarn}, anchor=west] at (-6.5, -2.0) {Couche Intégration, Automatisation \& APIs Externes};
        \node[comp] (e1) at (-4.2, -2.7) {Tâches Cron\\(ir.cron scheduler)};
        \node[comp] (e2) at (-1.3, -2.7) {API JSON-RPC\\/ Chatter Mail};
        \node[comp] (e3) at (1.8, -2.7) {API REST Google Gemini\\\& API Cerebras};
        \node[comp] (e4) at (5.0, -2.7) {Docker Engine\\\& WSL Runner};

        % Connecting Arrows
        \draw[<->, very thick, odooDark] (pres) -- (biz);
        \draw[<->, very thick, odooDark] (biz) -- (data);
        \draw[<->, very thick, odooDark] (biz) -- (ext);

    \end{tikzpicture}
    \end{adjustbox}
    \caption{Architecture logicielle multi-niveaux du module \odooModule}
    \label{fig:odoo_architecture}
\end{figure}

\subsection{Architecture de la Passerelle BYO-LLM et Sécurité des Outils}
Pour doter le progiciel de fonctionnalités d'assistance conversationnelle et de journalisation en langage naturel sans compromettre la sécurité, une passerelle \textbf{BYO-LLM} (\textit{Bring Your Own Large Language Model}) a été implémentée (\texttt{equipment\_llm\_providers.py}). 

Elle définit une classe de base abstraite \texttt{LLMProvider} exposant une interface normalisée (\texttt{generate()}, \texttt{chat()}, \texttt{structured\_output()}, \texttt{test\_connection()}) implémentée pour deux fournisseurs distincts :
\begin{enumerate}[leftmargin=1.2cm, label=\textbf{\alph*.}]
    \item \textbf{Google Gemini} (\texttt{GeminiProvider}) : utilise le modèle de dernière génération \texttt{gemini-2.5-flash} via l'API REST officielle Google Cloud, avec authentification par clé d'API protégée et support natif du schéma JSON strict pour les sorties structurées.
    \item \textbf{Cerebras} (\texttt{CerebrasProvider}) : s'interface avec l'API haute vitesse compatible OpenAI via jeton Bearer pour une exécution ultra-rapide.
\end{enumerate}

Le flux de traitement garantit l'isolation absolue des données d'entreprise :
\begin{equation}
\text{Entrée Utilisateur} \rightarrow \text{LLM} \rightarrow \text{Validation Schéma JSON} \rightarrow \text{Routage d'Outil Contrôlé} \rightarrow \text{Prévisualisation} \rightarrow \text{Confirmation} \rightarrow \text{Audit}
\end{equation}

Le modèle LLM ne dispose \textbf{d'aucun accès} à SQL, au shell du conteneur, à l'ORM ou au système de fichiers. Il génère uniquement des appels vers une table figée d'outils autorisés (\texttt{get\_equipment}, \texttt{create\_order}, \texttt{list\_orders}, \texttt{check\_parts}), dont l'exécution est soumise aux règles de sécurité Odoo (\textit{ir.model.access.csv}). Pour toute opération d'écriture, l'exécution est suspendue jusqu'à la confirmation explicite de l'utilisateur dans l'interface graphique.

\section{Inventaire Quantitatif du Code Source}
L'ensemble de la solution représente un volume de code hautement modulaire, résumé dans le tableau \ref{tab:code_inventory}.

\begin{table}[!ht]
\centering
\small
\renewcommand{\arraystretch}{1.15}
\begin{tabularx}{\textwidth}{l p{7.2cm} c r}
\toprule
\textbf{Composant} & \textbf{Description \& Rôle Fonctionnel} & \textbf{Fichiers} & \textbf{Lignes (LOC)} \\
\midrule
\textbf{Modèles Python} & 16 modèles métier (équipements, multi-sites, ordres, pièces, services IA, chat, logs) & 16 & 4\,120 \\
\textbf{Couche Compatibilité \& Providers} & Abstraction LLM Gemini/Cerebras (\texttt{equipment\_llm\_providers.py}) et compatibilité v17/v18 (\texttt{odoo\_compat.py}) & 2 & 780 \\
\textbf{Assistants \& Contrôleurs} & Wizards (Journalisation NL, Rapports d'exportation), Contrôleur QR Code & 4 & 490 \\
\textbf{Vues \& Menus XML} & Vues Kanban, Listes groupées, Formulaires, Calendriers, Vues Pivot/Graph, Menus & 14 & 2\,650 \\
\textbf{Données, Séquences \& Crons} & Séquences d'incrémentation, 18 tâches planifiées automatiques, données de base & 6 & 640 \\
\textbf{Sécurité \& Règles d'Accès} & 45 règles ACL (\textit{ir.model.access.csv}) et règles d'enregistrement au niveau ligne & 2 & 110 \\
\textbf{Rapports QWeb} & Déclarations de rapports et gabarits vectoriels d'impression PDF & 2 & 320 \\
\textbf{Tests Unitaires \& Mocks} & 138 tests backend Odoo (\texttt{tests/}) et 46 tests mock LLM (\texttt{test\_llm\_providers.py}) & 8 & 2\,480 \\
\textbf{Banc Playwright E2E} & 16 fichiers de spécifications de test de navigation navigateur (Chromium) & 16 & 1\,850 \\
\textbf{Scripts DevOps \& Orchestration} & Scripts d'exécution WSL/Docker, Dockerfiles et compose de test isolé & 6 & 280 \\
\midrule
\textbf{Total Général} & \textbf{Ensemble complet du système validé} & \textbf{76} & \textbf{13\,720} \\
\bottomrule
\end{tabularx}
\caption{Inventaire exhaustif des artefacts logiciels et métriques de code}
\label{tab:code_inventory}
\end{table}

\section{Spécification Détaillée des Modèles de Données ORM}
L'architecture de données repose sur 16 modèles Python étroitement interconnectés, présentés dans le tableau \ref{tab:models_spec}.

\begin{table}[!ht]
\centering
\small
\renewcommand{\arraystretch}{1.12}
\begin{tabularx}{\textwidth}{l l p{5.5cm} c}
\toprule
\textbf{Modèle ORM} & \textbf{Table PostgreSQL} & \textbf{Responsabilité Métier Principale} & \textbf{Champs} \\
\midrule
\texttt{equipment.equipment} & \texttt{equipment\_equipment} & Registre central, garanties, QR code, Health Score & 34 \\
\texttt{equipment.category} & \texttt{equipment\_category} & Familles et catégories hiérarchiques d'équipements & 12 \\
\texttt{equipment.location} & \texttt{equipment\_location} & Hiérarchie spatiale multi-sites (Site $\rightarrow$ Baie) & 14 \\
\texttt{equipment.maintenance.order} & \texttt{equipment\_maint\_order} & Ordre curatif/préventif, machine d'états, coûts & 30 \\
\texttt{repair.request} & \texttt{repair\_request} & Demandes d'intervention opérateurs et pannes urgentes & 18 \\
\texttt{equipment.checklist} & \texttt{equipment\_checklist} & Protocoles d'inspection et points de contrôle & 10 \\
\texttt{equipment.spare.part} & \texttt{equipment\_spare\_part} & Lignes de pièces consommées et liaison avec \texttt{stock} & 12 \\
\texttt{equipment.template} & \texttt{equipment\_template} & Modèles récurrents d'ordres de maintenance préventive & 16 \\
\texttt{equipment.dashboard} & \texttt{equipment\_dashboard} & Agrégation des KPIs et métriques en temps réel & 15 \\
\texttt{equipment.ai.config} & \texttt{equipment\_ai\_config} & Configuration des clés LLM, modèles et seuils d'activation & 14 \\
\texttt{equipment.ai.service} & \texttt{equipment\_ai\_service} & Algorithmes statistiques (MTBF, prévision budgétaire) & 12 \\
\texttt{equipment.ai.data} & \texttt{equipment\_ai\_data} & Extraction des vecteurs de features depuis l'historique & 8 \\
\texttt{equipment.ai.nl} & \texttt{equipment\_ai\_nl} & Parser en langage naturel (LLM + repli heuristique) & 11 \\
\texttt{equipment.ai.chat} & \texttt{equipment\_ai\_chat} & Assistant conversationnel interactif et routage d'outils & 15 \\
\texttt{equipment.ai.prediction}& \texttt{equipment\_ai\_pred} & Historique et piste d'audit inviolable des requêtes IA & 16 \\
\texttt{odoo.compat} & -- & Couche d'encapsulation de compatibilité Odoo 17/18 & -- \\
\bottomrule
\end{tabularx}
\caption{Spécification des 16 modèles de données de l'architecture ORM}
\label{tab:models_spec}
\end{table}

\subsection{Analyse Approfondie des Modèles Clés}

\paragraph{1. Modèle d'Emplacement Multi-Sites (\texttt{equipment.location})}
Ce modèle structure les actifs de l'entreprise au sein d'une arborescence parent-enfant stricte :
\begin{equation}
\text{Site} \rightarrow \text{Bâtiment} \rightarrow \text{Étage} \rightarrow \text{Salle/Atelier} \rightarrow \text{Baie/Emplacement}
\end{equation}
Chaque emplacement maintient un champ calculé \texttt{complete\_name} combinant les préfixes hiérarchiques et un compteur dynamique \texttt{equipment\_count} facilitant le filtrage dans l'interface utilisateur.

\paragraph{2. Modèle d'Ordre de Travail (\texttt{equipment.maintenance.order})}
Ce modèle pilote l'exécution opérationnelle. Il calcule automatiquement le coût total consolidé selon la formule :
\begin{equation}
\text{Coût Total} = \sum_{i} (Q_i \times C_i) + (T_{\text{int}} \times R_{\text{MO}})
\end{equation}
où $Q_i$ représente la quantité de la pièce de rechange $i$, $C_i$ son coût unitaire moyen pondéré, $T_{\text{int}}$ la durée d'intervention en heures et $R_{\text{MO}}$ le taux horaire de main-d'œuvre.

\paragraph{3. Modèle de Traçabilité et d'Audit IA (\texttt{equipment.ai.prediction})}
Conformément aux exigences suisses de protection des données (nLPD) et de conformité financière, toute sollicitation de la couche LLM (journalisation ou assistant) consigne un enregistrement d'audit horodaté contenant : le type d'opération, le texte utilisateur brut, la réponse structurée validée, les jetons consommés, le modèle utilisé et l'identifiant de confirmation de l'utilisateur.

\section{Algorithmique Avancée \& Qualification Réaliste de l'IA}

\subsection{Classification Réaliste des Composants d'Intelligence Artificielle}
Dans une démarche d'intégrité scientifique et technique, le tableau \ref{tab:ai_reality} établit la nature exacte de chaque composant d'intelligence artificielle intégré au module.

\begin{table}[!ht]
\centering
\small
\renewcommand{\arraystretch}{1.2}
\begin{tabularx}{\textwidth}{l p{3.2cm} c p{6.0cm}}
\toprule
\textbf{Fonctionnalité} & \textbf{Classification} & \textbf{ML Entraîné ?} & \textbf{Mécanisme \& Preuve Logicielle} \\
\midrule
\textbf{Health Score} & Basé sur des Règles & Non & Pondération déterministe multicritère sur 100 points \\
\textbf{Risque de Panne} & Statistique & Non & Score pondéré à 6 facteurs basé sur MTBF/MTTR \\
\textbf{Planification Intelligente} & Basé sur des Règles & Non & Déclencheurs sur seuils d'échéance et compteurs \\
\textbf{Prévision des Coûts} & Statistique & Non & Moyennes mobiles pondérées avec correction saisonnière \\
\textbf{Conseil de Remplacement} & Statistique & Non & Analyse d'amortissement et retour sur investissement \\
\textbf{Dépréciation d'Actif} & Statistique & Non & Calcul d'amortissement linéaire selon durée de vie \\
\textbf{Journalisation NL} & \textbf{LLM Génératif} & -- & Appels réels Google Gemini avec repli par mots-clés \\
\textbf{Assistant Conversationnel} & \textbf{LLM Génératif} & -- & Routage d'outils et validation structurée Gemini \\
\textbf{Détection d'Anomalies} & Non Implémenté & -- & Exclu du périmètre fonctionnel par spécification \\
\bottomrule
\end{tabularx}
\caption{Classification rigoureuse des fonctionnalités d'intelligence artificielle}
\label{tab:ai_reality}
\end{table}

\subsection{Algorithme de Score de Santé (\textit{Health Score})}
Le score de santé $H \in [0, 100]$ fournit une métrique instantanée de la fiabilité de chaque machine :
\begin{equation}
H = 0.20 \cdot S_{\text{Age}} + 0.25 \cdot S_{\text{Pannes}} + 0.20 \cdot S_{\text{Downtime}} + 0.25 \cdot S_{\text{Adherence}} + 0.10 \cdot S_{\text{Garantie}}
\end{equation}
Chaque sous-score est borné dans $[0, 100]$ avec protection native contre les divisions par zéro lors de l'initialisation de machines neuves n'ayant aucun historique de fonctionnement.

\subsection{Prédiction Statistique de Défaillance (MTBF \& MTTR)}
La probabilité de panne sous 30 jours repose sur l'estimation du Temps Moyen Entre Pannes (\textit{Mean Time Between Failures - MTBF}) et du Temps Moyen de Réparation (\textit{Mean Time To Repair - MTTR}) :
\begin{equation}
\text{MTBF} = \frac{\text{Heures Totales de Fonctionnement}}{\max(1, N_{\text{pannes}})}, \quad \text{MTTR} = \frac{\text{Heures Totales d'Arrêt}}{\max(1, N_{\text{pannes}})}
\end{equation}
Un score de risque composite à 6 composantes pondérées signale les équipements prioritaires nécessitant une révision immédiate.

\section{Validation de Navigation et Tests E2E sous Playwright}

\subsection{Mise en Place du Banc de Test Navigateur E2E}
Pour combler l'écart critique entre les tests unitaires de backend (qui n'exécutent pas de moteur de rendu JavaScript/OWL) et l'expérience utilisateur réelle, un banc complet de tests \textit{End-to-End} (E2E) a été déployé à l'aide de \textbf{Playwright} et du navigateur \textbf{Chromium}.

L'orchestration repose sur le script d'automatisation \texttt{scripts/e2e-run.sh} exécuté sous environnement WSL/Docker, garantissant la persistance du conteneur Odoo 18, la surveillance active de sa disponibilité (\textit{health polling}) et la capture d'artéfacts de test (traces de débogage, captures d'écran et vidéos de session WebM).

\subsection{Couverture Complète de la Navigation Inter-Vues}
La campagne E2E comprend \textbf{16 suites de spécifications} exécutant \textbf{32 tests complets} (taux de réussite de 100\%, zéro échec, zéro test ignoré). Elle valide minutieusement l'ensemble des parcours de navigation au sein de l'interface Odoo 18 :

\begin{enumerate}[leftmargin=1.2cm, label=\textbf{\arabic*.}]
    \item \textbf{Navigation en Cartes Kanban et Vues Listes Groupées} : Odoo 18 introduit un comportement par défaut où les listes groupées s'affichent à l'état replié (\textit{collapsed}). Les tests E2E ont validé le mécanisme d'expansion automatique des groupes et la sélection dynamique des cartes d'équipements dotées de badges de couleur de santé.
    \item \textbf{Navigation dans l'Arborescence Multi-Sites} : Parcours hiérarchique depuis le site parent jusqu'aux baies individuelles, validant le filtrage instantané des équipements hébergés.
    \item \textbf{Navigation Formulaire et Fil d'Ariane (\textit{Breadcrumbs})} : Vérification de la cohérence de navigation lors du passage d'une vue liste vers une vue fiche, validation de l'avancement sur la barre d'état (\textit{Draft $\rightarrow$ Confirmed $\rightarrow$ In Progress $\rightarrow$ Done}) et retour via le fil d'Ariane.
    \item \textbf{Navigation Calendrier et Planning Adapté} : Remplacement de la vue propriétaire \texttt{gantt} (qui déclenchait une erreur 500 sur Odoo 18 Community) par une vue de planification calendrier mensuelle et hebdomadaire entièrement validée sous Playwright.
    \item \textbf{Génération et Navigation QR Code Mobile} : Contrôle du rendu vectoriel du QR code, validation de l'intégrité de l'entête binaire PNG généré et affichage sans latence de la fiche mobile associée.
    \item \textbf{Navigation Interactive du Tableau de Bord} : Validation du clic sur les widgets KPI (ordres en retard, machines critiques) déclenchant le drilldown vers les listes pré-filtrées correspondantes.
    \item \textbf{Navigation des Dialogues Modaux de Journalisation NL} : Résolution du comportement où le wizard se fermait immédiatement après analyse ; validation du rechargement modal permettant d'afficher les résultats extraits avant la création de l'ordre.
    \item \textbf{Navigation de l'Assistant Conversationnel} : Maintien de l'état de conversation multi-tours dans le dialogue modal lors des échanges et confirmations.
    \item \textbf{Validation du Flux de Confirmation d'Écriture} : Preuve formelle par les tests E2E qu'aucune écriture en base n'est exécutée lors d'une requête de modification tant que l'utilisateur n'a pas cliqué sur le bouton de confirmation dans l'interface.
    \item \textbf{Scénario Complet de Démonstration (Full Demo Scenario)} : Exécution continue et sans faille de la chaîne complète de maintenance en 20 étapes successives en \textbf{1,1 minute} (\texttt{tests/e2e/full-demo.spec.ts}).
\end{enumerate}

\subsection{Anomalies Réelles Détectées et Corrigées par la Validation de Navigation}
La valeur ajoutée majeure de la campagne de tests Playwright réside dans la détection et la correction de \textbf{10 anomalies réelles d'interface et de navigation} qui échappaient totalement à la suite de tests backend, synthétisées dans le tableau \ref{tab:e2e_bugs}.

\begin{table}[!ht]
\centering
\small
\renewcommand{\arraystretch}{1.18}
\begin{tabularx}{\textwidth}{c p{3.6cm} p{4.2cm} X}
\toprule
\textbf{N$^\circ$} & \textbf{Composant / Vue} & \textbf{Comportement Défaillant Détecté} & \textbf{Solution Technique Appliquée} \\
\midrule
1 & Liste groupée équipements & Groupes repliés masquant les lignes en vue liste Odoo 18 & Ajout d'une séquence explicite de dépliage de groupe dans le framework d'UI \\
2 & Action Ordres de Travail & Erreur 500 (« Oops ») au chargement de l'action & Retrait de \texttt{gantt} du \texttt{view\_mode} (vue exclusive à Odoo Enterprise) \\
3 & Lignes pièces détachées & L'administrateur voyait 0 ligne saisie par les techniciens & Correction de la règle d'enregistrement avec \texttt{perm\_read="0"} \\
4 & Dialogue Assistant NL & Le modal de saisie se fermait avant d'afficher les résultats & Retour d'une action de rechargement (\textit{reload action}) dans \texttt{action\_parse} \\
5 & Chat Assistant conversationnel & Le dialogue se fermait lors de l'envoi ou confirmation & Application d'actions de rechargement sur \texttt{action\_send} et \texttt{confirm} \\
6 & Détection d'équipements & Questions directes renvoyaient l'aide par défaut & Implémentation d'une recherche par sous-chaîne sur noms, codes et séries \\
7 & Générateur de QR Code & Rejet injustifié des codes QR de taille inférieure à un seuil & Validation de la signature binaire PNG au lieu d'un seuil de taille arbitraire \\
8 & Données de démonstration & Échec de chargement des produits consommables & Migration vers \texttt{'consu'} avec \texttt{is\_storable=True} conforme v18 \\
9 & Sélecteurs de test E2E & Évaluation silencieuse erronée des tuples JS & Correction des domaines Playwright sous forme de tableaux réels \texttt{[["f","=",x]]} \\
10 & Génération par gabarit & Conflit d'assertion sur la génération globale d'ordres & Restriction du périmètre de vérification aux ordres issus du gabarit testé \\
\bottomrule
\end{tabularx}
\caption{Synthèse des 10 anomalies d'ergonomie et de navigation résolues par Playwright}
\label{tab:e2e_bugs}
\end{table}

\section{Télémétrie d'Assurance Qualité et Preuves sur Disque}
La validation de la solution repose sur une batterie complète de tests automatisés couvrant tous les niveaux de la pile logicielle. L'ensemble des résultats est attesté par des fichiers de logs et des enregistrements vidéo sur disque, synthétisés dans le tableau \ref{tab:qa_telemetry}.

\begin{table}[!ht]
\centering
\small
\renewcommand{\arraystretch}{1.2}
\begin{tabularx}{\textwidth}{l c c p{4.2cm}}
\toprule
\textbf{Batterie de Tests} & \textbf{Volume} & \textbf{Résultat} & \textbf{Preuve Matérielle sur Disque} \\
\midrule
\textbf{Tests Backend Odoo 18} & 138 tests & \textbf{138/138 PASS (100\%)} & \texttt{equipment\_maintenance/.bt\_final.log} (RC=0) \\
\textbf{Tests Mock Providers LLM} & 46 tests & \textbf{46/46 PASS (100\%)} & \texttt{tests/test\_llm\_providers.py} (unittest) \\
\textbf{Tests Live API Gemini} & 6 tests & \textbf{6/6 PASS (100\%)} & \texttt{tests/live\_api\_results.json} (\texttt{gemini-2.5-flash}) \\
\textbf{Tests Live API Cerebras} & 2 tests & \textbf{Auth PASS / Quota 402} & Requête HTTP réelle validée ; génération bloquée quota \\
\textbf{Tests E2E Navigation Playwright} & 32 tests & \textbf{32/32 PASS (100\%)} & \texttt{artifacts/playwright-report/index.html} \\
\textbf{Scénario Démo Système Complet} & 20 étapes & \textbf{PASS (1,1 minute)} & \texttt{artifacts/videos/full-system-demo.webm} (2,9 Mo) \\
\textbf{Scénario Démo IA / BYO-LLM} & 5 étapes & \textbf{PASS} & \texttt{artifacts/videos/ai-llm-demo.webm} (1,3 Mo) \\
\bottomrule
\end{tabularx}
\caption{Matrice consolidée de télémétrie d'assurance qualité logicielle}
\label{tab:qa_telemetry}
\end{table}

\section{Conclusion}
Ce chapitre a exposé l'architecture complète, l'implémentation des 16 modèles ORM et la concrétisation de la passerelle d'intelligence artificielle BYO-LLM. L'apport fondamental de la campagne de tests Playwright a permis de valider la navigation fluide à travers toutes les vues d'Odoo 18 et d'éliminer 10 anomalies critiques d'expérience utilisateur. Avec 138 tests backend, 46 tests mock LLM et 32 tests E2E au vert, le module \textbf{\odooModule} démontre une robustesse de classe industrielle, prêt pour un déploiement en production chez les clients de \textbf{\companyName}.
